% Sección 11 - Verificación de los relojes de Lamport (Dayver)
% Borra los comentarios cuando escribas el texto.

\subsection{Cómo se verificó}
% DATOS: el script analysis/verificar_lamport.py lee los 234 CSV de eventos
% (data/eventos_<mecanismo>_<nodo>_<equipo>_<corrida>.csv) de las 78 corridas.
% Empareja cada recepción con su envío: una recepción guarda en reloj_emisor el
% timestamp que traía el mensaje, así que su envío es la fila con
%   (corrida, nodo, reloj) del envío = (corrida, origen, reloj_emisor) de la recepción.
% Dentro de un nodo el reloj nunca se repite, por eso esa clave es única.
% Revisa: que toda recepción tenga su envío, que reloj de la recepción > reloj
% del envío, que todo envío tenga su recepción y que en cada nodo el reloj suba
% en cada fila. Verifica los datos reales: no genera valores.

\subsection{Resultado}
% DATOS (de analysis/lamport_resumen.csv y lamport_violaciones.csv):
% 20 400 recepciones en 78 corridas; las 20 400 cumplen reloj_recepcion >
% reloj_envio (100 %). Cero recepciones sin envío, cero envíos sin recepción,
% cero violaciones (lamport_violaciones.csv solo tiene el encabezado).
% Ojo: son 20 400 porque cuenta también las corridas de latencia (los nodos
% usan el reloj en todos los modos), no solo las 30 de intercambios.

\subsection{Dos eventos concurrentes con el mismo reloj}
% DATOS: corrida n3-r1 de sockets (CSV data/eventos_sockets_nodo-A_mcastro_n3-r1.csv
% y data/eventos_sockets_nodo-B_mcastro_n3-r1.csv):
%  - nodo-A: envío a nodo-B con reloj 1 (hora_real 1795365707.206059)
%  - nodo-B: envío a nodo-C con reloj 1 (hora_real 1795365707.207598)
%  Son los primeros eventos de cada nodo: ninguno había recibido nada, así que
%  ninguno pudo influir en el otro -> son concurrentes (a || b).
%  Por la hora real, A envió 1,5 ms antes que B, pero el reloj no lo refleja.
% POR QUÉ LAMPORT NO LOS DISTINGUE (explícalo tú):
%  - Lamport garantiza: si a -> b (a causa b), entonces C(a) < C(b).
%  - Lo contrario no vale: C(a) < C(b) no implica a -> b, y C(a) = C(b) no
%    dice cuál fue primero. Con el mismo valor solo se sabe que ninguno causó
%    al otro (si uno causara al otro, su reloj sería menor).
%  - Es un solo número por evento: no guarda de qué nodos viene la información.
%    Para distinguir concurrencia de causalidad se necesitan relojes vectoriales.
%  - Para un orden total se desempata con el nombre del nodo (A antes que B),
%    pero ese orden es una convención, no lo que pasó.

\subsection{Convergencia}
% DATOS: no repitas la tabla. Refiere la de la sección 7: \ref{tab:convergencia}.
% Las 30 corridas convergen; 2, 6 y 8 intercambios con 2, 3 y 4 nodos, con
% reloj 2, 6 y 8 en todas las repeticiones, porque cada nodo recibe de un solo
% antecesor por una sola conexión y el orden de los eventos siempre es el mismo.
% El tiempo sí varía; la explicación está en \ref{sec:discusion-lamport}.
